\documentclass[10pt,a4paper]{amsart}
\usepackage[margin=19mm]{geometry}
\usepackage{amsmath,amssymb,mathtools}
\usepackage[scr=rsfs,cal=euler]{mathalfa}
\usepackage{xfrac}
\usepackage[unicode,hidelinks,bookmarks=false]{hyperref}
\newcommand{\ie}{i.e.\ }
\newcommand{\cf}{cf.\ }
\newcommand{\resp}{respectively\ }
\newcommand{\Subsection}{\subsection}
\providecommand{\nobreakdash}{\nobreak\hbox{-}}
\setlength{\emergencystretch}{2em}
\multlinegap0pt
\usepackage{bm}
\usepackage{bbm}
\usepackage{dsfont}
\usepackage{xspace}
\usepackage{cancel}
\usepackage{mathtools}
\usepackage{amsthm}
\makeatletter
\@ifundefined{defi}{\newtheorem{defi}{Defi}}{}
\@ifundefined{lem}{\newtheorem{lem}[defi]{Lemma}}{}
\@ifundefined{propal}{\newtheorem{propal}[defi]{Proposition}}{}
\makeatother
\title[High-order long-time asymptotics for small solutions to the ]{High-order long-time asymptotics for small solutions to the one-dimensional nonlinear Schr\"odinger equation}
\author{Jacek Jendrej, Tony Salvi}
\date{Source: arXiv:2602.19374v1 (CC BY 4.0)}
\begin{document}
\maketitle

\noindent\emph{Source and licence.} High-order long-time asymptotics for small solutions to the one-dimensional nonlinear Schr\"odinger equation. Version arXiv:2602.19374v1 (CC BY 4.0), distributed under the Creative Commons Attribution 4.0 International licence, as recorded in the arXiv licence metadata for this exact version and shown on the abstract page https://arxiv.org/abs/2602.19374v1. Full rights notice: licenses/arxiv-2602.19374-CC-BY-4.0.txt.\par\smallskip

\noindent\textit{Editorial scope.} Counted unit: the Propal whose source label is \texttt{propal:controlerror} in the article (source0.tex lines 460--471, source label \texttt{propal:controlerror}), delivered together with the article's own complete original proof of it (source0.tex lines 565--622, one \texttt{proof} environment of 58 lines). No mathematics is written, repaired or re-derived: statement and proof are the article's own text line for line. The proof is detached from the statement in the source: the article states the result at lines 460--471 and prints its proof at lines 565--622, and the proof environment's own header names the statement, so the association is the article's own. Required context retained uncounted: lem:Rtrans (lines 473--484); lem:Rcalculus (lines 529--536). Every definition, display and label of those lines is retained unchanged. Omitted from this package and not redistributed: 1--459, 472--472, 485--528, 537--564, 623--1536. Nothing inside a retained excerpt is omitted or altered: every retained body is delivered exactly as the article's lines are written, so no omission receipt of any kind accompanies this package. Citations: the retained excerpts carry no citation command at all, so no bibliography entry of the article is transcribed and no reference is redistributed. Reference adaptation: none was needed and none was made; every reference inside the delivered unit resolves inside it, so no unresolved reference or citation is produced. Printed slips kept literal: no printed word, symbol, hypothesis or number of the counted statement or of its proof was altered. Layout and class: the article's own class is \texttt{article} with packages including inputenc, fontenc, babel, tikz, geometry, graphicx, amssymb, biblatex, lineno, amsthm, amsmath, amsfonts, mathtools, latexsym; the wrapper is the lane's pinned \texttt{amsart} editorial wrapper at 10pt on A4, which restates the theorem-like environments the retained text needs and adds the article's own macro definitions that the retained text needs (none), with extra wrapper packages (bm, bbm, dsfont, xspace, cancel, mathtools). The delivered numbering is the unit's own. Assets: no retained span contains a figure environment, a graphics-inclusion command, a PDF-page inclusion, a table or a TikZ picture; no image file is copied into this package, the package declares no asset dependency and no image file is redistributed with it. Licence: version arXiv:2602.19374v1 of this article is distributed under the Creative Commons Attribution 4.0 International licence, as recorded in the arXiv licence metadata for this exact version and shown on the abstract page https://arxiv.org/abs/2602.19374v1. A full rights notice is delivered at licenses/arxiv-2602.19374-CC-BY-4.0.txt. Characters: every retained excerpt is pure ASCII, so the delivered engine, XeLaTeX with its default Latin Modern text font, reports no missing character.
\begin{lem}
\label{lem:Rtrans}
For $s\in[1,+\infty)$, $\xi\in\mathbb{R}$ and $\sigma\in\mathbb{R}^{2n+1}$, we have 
\begin{align}
\label{equation:Rtransform}
    \mathcal{F}^{-1}_{\eta}[\mathcal{N}^n](s,\xi,\sigma)=\frac{1}{(2\pi)^{1/2}}e^{i\xi(\sum^{2n}_{j=1}\sigma_{j})}\int e^{-i\xi x}\mathcal{M}^n(s,x,\sigma)dx
\end{align}
with 
\begin{align*}
   \mathcal{M}^n=\left(\prod^{n}_{j=1}f(x-\sigma_{2j}-\sum^{2(j-1)}_{k=1}\sigma_k-\sum^{n-1}_{k=j}\sigma_{2k+1})\overline{f(x-\sum^{2(j-1)}_{k=1}\sigma_{k}-\sum^{n-1}_{k=j}\sigma_{2k+1})}\right)f(x-\sum^{n-1}_{k=0}\sigma_{2k+1}).
\end{align*}
\end{lem}

\begin{lem}
\label{lem:Rcalculus}
For $\mathcal{M}^n$ defined as in the previous lemma \ref{lem:Rtrans} and some $\alpha\geq0$, we have 
\begin{align}
\label{equation:Rcalculus}
\int\int (\sum_{j=1}^{2n}|\sigma_j|+|x|)^{\alpha}\mathcal{M}^n(x,\sigma)dxd\sigma\lesssim\int (\sum_{j=1}^{2n+1}|Y_{j}|)^{\alpha}\left(\prod^{n}_{j=1}|f(Y_{2j+1})||\overline{f(Y_{2j})}|\right)|f(Y_{1})|dY_1\dots dY_{2n+1}
\end{align}
\end{lem}

\begin{propal}
\label{propal:controlerror}
Given $p\in\mathbb{N}$. If $||t^{-\alpha}\langle x\rangle^{p+1} f||_{L^\infty_tL^2_x}\leq C$ for some $1/(16d+8)>\alpha>0$ and some $C>0$ then, for $r,q\in\mathbb{N}$ such that $2r+q-2=p$ if $r>0$ or $p=q$ if $r=0$, there exists $1>\delta>0$ such that for $s\geq1$
\begin{align}
&||\mathfrak{R}^n_r(s)||_{W^{q,\infty}_\xi}\lesssim s^{-n-r+1-\delta}C^{2n+1}.
\end{align}
Moreover, if $||\hat{w}||_{L^\infty_tW^{p,\infty}_\xi}\leq C$, then, under the same relations for $q,r$ and $p$, there exists $1>\delta>0$ such that for $s\geq1$
\begin{align}
&||\Theta\mathfrak{R}^n_r(s)||_{W^{q,\infty}_\xi}\lesssim s^{-n-r+1-\delta}h(C)C^{2n+1}
\end{align}
for some polynomial $h$.
\end{propal}

\begin{proof}[Proof of proposition \ref{propal:controlerror}]
For $r>0$, we have
\begin{align*}
    \partial^q_\xi \mathfrak{R}^n_r(s)&=\lambda_n\frac{1}{(2\pi)^n}\int\frac{1}{s^n}(e^{-i(\eta\cdot  (Q^n)^{-1}\eta)/(2s)}-\sum^{r-1}_{j=0}\frac{(-i\eta\cdot (Q^n)^{-1}\eta)^j}{(2s)^{j}j!}) \partial^q_\xi\mathcal{F}^{-1}_{\eta}[\mathcal{N}^n](s,\xi,\eta)d\eta\\
    &\lesssim \int\frac{1}{s^{n+r-\beta}}(\eta\cdot (Q^n)^{-1}\eta)^{r-\beta} |\partial^q_\xi\mathcal{F}^{-1}_{\eta}[\mathcal{N}^n](s,\xi,\eta)|d\eta\\
    &\lesssim\int\int\frac{1}{s^{n+r-\beta}}|\eta|^{2r-2\beta} (\sum_{j=1}^{2n}|\eta_j|+|x|)^{q}|\mathcal{M}^{n}(s,x,\eta)|dxd\eta\\
    &\lesssim\frac{1}{s^{n+r-\beta}}\int\int(\sum_{j=1}^{2n}|\eta_j|+|x|)^{2r+q-2\beta} |\mathcal{M}^{n}(s,x,\eta)|dxd\eta
\end{align*}
with $0<\beta<1$, and for $r=0$
\begin{align*}
    \partial^q_\xi \mathfrak{R}^n_0(s)&=\lambda_n\frac{1}{(2\pi)^n}\int\frac{1}{s^n}(e^{-i(\eta\cdot  (Q^n)^{-1}\eta)/(2s)}) \partial^q_\xi\mathcal{F}^{-1}_{\eta}[\mathcal{N}^n](s,\xi,\eta)d\eta\\
    &\lesssim \frac{1}{s^{n}} \int\int(\sum_{j=1}^{2n}|\eta_j|+|x|)^{q}|\mathcal{M}^{n}(s,x,\eta)|dxd\eta.
\end{align*} 
where we use lemma \ref{lem:Rtrans} to compute $\mathcal{F}^{-1}_{\eta}[\mathcal{N}^n](s,\xi,\eta)$ in both cases. Then, applying lemma \ref{lem:Rcalculus} leads to
\begin{align*}
    \partial^q_\xi \mathfrak{R}^n_r(s)\lesssim\frac{1}{s^{n+r-\beta}}\int\int(\sum_{j=1}^{2n+1}|Y_{j}|)^{2r+q-2\beta}\left(\prod^{n}_{j=1}|f(s,Y_{2j+1})||\overline{f(s,Y_{2j})}|\right)|f(s,Y_1)|dY_1\dots dY_{2n+1}
\end{align*}
and 
\begin{align*}
    \partial^q_\xi \mathfrak{R}^n_0(s)\lesssim\frac{1}{s^{n}} \int\int(\sum_{j=1}^{2n+1}|Y_{j}|)^{q}\left(\prod^{n}_{j=1}|f(s,Y_{2j+1})||\overline{f(s,Y_{2j})}|\right)|f(s,Y_1)|dY_1\dots dY_{2n+1}.
\end{align*}
This gives $2n+1$ integrals with weighted $f$'s as integrands that can be computed one by one. For $r>0$, we get 
\begin{align*}
    \partial^q_\xi \mathfrak{R}^n_r(s)
    &\lesssim \frac{1}{s^{n+r-\beta}}\sum_{\overrightarrow{m}\in\mathbb{N}^{2n+1}, |\overrightarrow{m}|=2r+q}\prod^{2n+1}_{j=1}||\langle x\rangle^{m_j\gamma} f(s)||_{L^1_x} \\
    &\lesssim \frac{1}{s^{n+r-\beta-(2n+1)\alpha}}\sum_{\overrightarrow{m}\in\mathbb{N}^{2n+1}, |\overrightarrow{m}|=2r+q}\prod^{2n+1}_{j=1}||s^{-\alpha}\langle x\rangle^{m_j\gamma+1/2+\varepsilon} f(s)||_{L^2_x} \\
    &\lesssim \frac{1}{s^{n+r-\beta-(2n+1)\alpha}}C^{2n+1}
\end{align*}
for $|\overrightarrow{m}|=\sum_{j=1}^{2n+1}|\overrightarrow{m}_j|$ and $\beta>3/4$, and where $0<\gamma<1$ is such that $2r+q-2\beta=(2r+q)\gamma$ and $\varepsilon>0$ such that $2\beta-3/2>\varepsilon$. With a simpler reasoning, we have 
\begin{align*}
    \partial^q_\xi \mathfrak{R}^n_0(s)
    &\lesssim \frac{1}{s^{n-(2n+1)\alpha}}\sum_{\overrightarrow{m}\in\mathbb{N}^{2n+1}, |\overrightarrow{m}|=q}\prod^{2n+1}_{j=1}||s^{-\alpha}\langle x\rangle^{m_j+1/2+\varepsilon} f(s)||_{L^2_x} \\
    &\lesssim \frac{1}{s^{n-(2n+1)\alpha}}C^{2n+1}
\end{align*}
for $1/2>\varepsilon>0$. Thus, if $\alpha<1/(16d+8)$ and $\beta>3/4$ is small enough, we obtain for some $1>\delta>0$
\begin{align*}
    \partial^q_\xi \mathfrak{R}^n_r(s)\lesssim s^{-n-r+1-\delta}C^{2n+1}
\end{align*}
for $s\geq 1$.
Then, bounds on lower derivatives can be obtain the same way:
\begin{align}
|| \mathfrak{R}^n_r(s)||_{W^{q,+\infty}_x}\lesssim s^{-n-r+1-\delta}C^{2n+1}.
\end{align}
This is the first statement of the proposition. Now, knowing that 
\begin{align*}
    \Theta(s,\xi)=e^{i\int^s_1\frac{|\hat{f}(\tau,\xi)|^2}{\tau}d\tau}
\end{align*}
and that $|\hat{f}(\tau,\xi)|=|\hat{w}(\tau,\xi)|$, we get that, for $k\leq q$,
\begin{align*}
    ||\partial^k_\xi \Theta(s)||_{L^\infty_x}&\lesssim \sum^{k}_{j=1}\ln(s)^j||\hat{w}||^{2j}_{L^\infty_t W^{k,\infty}_x}\\
   & \lesssim \sum^{k}_{j=1}\ln(s)^jC^{2j}.
\end{align*}
Finally, up to a change of $1>\delta>0$, 
\begin{align}
||\Theta\mathfrak{R}^n_r(s)||_{W^{q,+\infty}_x}\lesssim s^{-n-r+1-\delta}h(C)C^{2n+1}
\end{align}
for $s\geq 1$ and some polynomial $h$. 
\end{proof}

\end{document}
